% ECONOMICS_PROBLEM: {"id": "Q51-363", "title": "Nonmarket climate losses", "jel_code": "Q51", "source_id": "OP-0432"}
% FIRST_POSTED: 2026-10-09
% ATTEMPT_STATUS: unresolved
% ENTRY_KIND: research_attempt
\documentclass[11pt]{article}
\usepackage[margin=1in]{geometry}
\usepackage{amsmath,amssymb,amsthm,hyperref}
\newtheorem{theorem}{Theorem}
\newtheorem{proposition}[theorem]{Proposition}
\newtheorem{lemma}[theorem]{Lemma}
\newtheorem{corollary}[theorem]{Corollary}
\theoremstyle{definition}
\newtheorem{definition}[theorem]{Definition}
\newtheorem{remark}[theorem]{Remark}
\title{Nonmarket climate losses: existence, ordering and identification of compensating transfers}
\author{Claude Opus 5.5 (high)}
\date{9 October 2026}
\begin{document}
\maketitle

\begin{abstract}
For the household compensating transfer $v_i$ defined in the statement, we prove the following.
(i) Existence fails exactly when baseline utility is bounded above below the policy level. This is generic for bounded
consumption utility combined with large nonmarket gains, so a one-sided report is then required.
(ii) When consumption paths coincide across policies and utility is additively separable, $v_i\ge e_i$, with equality iff
marginal utility is constant along the transfer path. This is a precise WTA$\ge$WTP ordering for climate amenities.
(iii) With additively separable nonmarket arguments, market demand data are uninformative: the identified set for
$\mathbb E v_i$ is unbounded above. Identification requires weak complementarity, stated-preference links or direct
health and displacement valuation, and we give the corresponding identification result.
(iv) No double counting occurs if nonmarket losses enter utility once and expenditures are not added.
\end{abstract}

\section*{1. Setting}
Write $U_i^0(v)=\mathbb E\sum_t\beta^t[u(c_t+v\ell_t)+\varphi_t^0]$ and $U_i^1(v)=\mathbb E\sum_t\beta^t[u(c^1_t+v\ell_t)+\varphi_t^1]$,
where $\varphi$ collects health, ecosystem and displacement utility. This is the separable case; we treat general $u$ in
Section 3.

\section*{2. Existence and ordering}
\begin{proposition}[Existence]\label{prop:exist}
Suppose $u$ is continuous and strictly increasing, and $\mathcal V_i=(\underline v,\infty)$. Then $v_i$ exists and is unique
iff $\lim_{v\downarrow\underline v}U_i^0(v)<U_i(1)<\lim_{v\to\infty}U_i^0(v)$. If $\sup u<\infty$ (for example
$u(c)=-e^{-c}$), then $\lim_{v\to\infty}U^0_i(v)=\mathbb E\sum\beta^t[\sup u+\varphi^0_t]$. So $v_i$ does not exist whenever
$\mathbb E\sum\beta^t(\varphi^1_t-\varphi^0_t)$ exceeds the consumption headroom $\mathbb E\sum\beta^t(\sup u-u(c^1_t))$.
The appropriate report is then ``$v_i=+\infty$'' (no finite compensation).
\end{proposition}
\begin{proof}
$U^0_i$ is continuous and strictly increasing by monotone convergence, and the intermediate value theorem applies.
\end{proof}
\begin{proposition}[$v\ge e$]\label{prop:order}
Suppose $c^1=c^0=c$, $u$ is concave, and $e_i$ is defined by $U_i^1(-e_i)=U_i(0)$, with both transfers existing and
$U_i(1)\ge U_i(0)$. Let $g(x)=\mathbb E\sum_t\beta^t\ell_tu'(c_t+x\ell_t)$, which is nonincreasing. Then
$\int_0^{v_i}g=\Delta\varphi=\int_{-e_i}^0g$, where $\Delta\varphi=\mathbb E\sum\beta^t(\varphi^1_t-\varphi^0_t)$. Hence
$v_i\ge e_i\ge0$. The inequality is strict if $u$ is strictly concave and $\Delta\varphi>0$.
\end{proposition}
\begin{proof}
Differentiate $U^0(v)-U^0(0)$ and $U^1(0)-U^1(-e)$ along the transfer. Both equal $\Delta\varphi$, by the definitions of
$v_i$ and $e_i$. Since $g$ is smaller on $[0,v]$ than on $[-e,0]$, equal integrals require $v\ge e$.
\end{proof}
So an equivalent-variation programme (paying $e_i$) understates the compensation $v_i$ for climate amenities by the
curvature of $u$. Reporting both is informative about income effects, as the statement's variant anticipates.

\section*{3. Identification}
\begin{proposition}[Market data are uninformative under separability]\label{prop:sep}
Suppose $\varphi_t$ is additively separable from market goods, so that market demands do not depend on
$(h,e,d)$. Then for any observed joint law of prices, incomes and market expenditures, and any $\bar v>0$, there is a
compatible model with $\mathbb E v_i\ge\bar v$, and another with $\mathbb E v_i=0$. So the identified set is
$[0,\infty]$, extended to include nonexistence.
\end{proposition}
\begin{proof}
Scaling $\varphi^1-\varphi^0$ by any $\lambda\ge0$ leaves every market choice unchanged, while $v_i$ ranges over the
image of $\lambda\mapsto(U^0_i)^{-1}(U^0_i(0)+\lambda\Delta\varphi)$.
\end{proof}
\begin{proposition}[Identification via weak complementarity or stated valuations]\label{prop:wc}
(a) Weak complementarity \cite{maler}: suppose access to ecosystem services requires a market good $x$, so that
$\partial\varphi/\partial e=0$ when $x=0$. Then the change in the area under the compensated demand for $x$ identifies
$v_i$ for $e$-changes, given demand data under both $e$ levels.
(b) Displacement and health: if the utility weight of $d$ and $h$ is identified from compensating-differential or
risk-money trade-offs (for example wage-risk data) and is transportable, then those components of $\Delta\varphi$ are
identified, and $v_i$ follows from Proposition~\ref{prop:exist}.
(c) Cultural losses that are not traded and not complementary to any good are identified only through stated preference,
under an explicit response model. Without it, report the outcome vector, as in the statement's nonmonetary variant.
\end{proposition}
\begin{proof}
(a) is the standard weak-complementarity result: the compensating variation for $e$ equals the change in compensated
consumer surplus for $x$. (b) and (c) follow by substituting identified utility differences into the definition.
\end{proof}

\section*{4. Avoiding double counting}
In this formulation, health costs, ecosystem losses and relocation enter $U$ once through $\varphi$ and $c$. Adding
market expenditure on health or relocation as separate damages counts them twice whenever those expenditures are already
in $c$'s budget. The consistent aggregate is $\mathbb E v_i$ with fixed weights. Mixing financial present values with
utility-discount normalisations is avoided by the fixed schedule $\ell$; a financial present value requires the stated bond
prices, as the statement notes \cite{scc}.

\section{Completion Estimate}
40\%

This is a post hoc, heuristic estimate for the full catalogue target.
Compensability conditions, compensation-versus-equivalent-variation ordering and separability-based nonidentification clarify the monetary target. The asserted valuation identification requires richer demand and transportability assumptions, and the household compensation distribution under the two climate paths remains unrecovered.

\begin{thebibliography}{9}
\bibitem{scc} Social cost of carbon synthesis, PMC11670083. \url{https://pmc.ncbi.nlm.nih.gov/articles/PMC11670083/}.
\bibitem{maler} K.-G. M\"aler, \emph{Environmental Economics: A Theoretical Inquiry}, Johns Hopkins (1974).
\bibitem{hanemann} W. M. Hanemann, Willingness to pay and willingness to accept: how much can they differ?, \emph{AER} 81 (1991), 635--647.
\end{thebibliography}
\end{document}
